\documentstyle[% 


righttag% 


,11pt% 


,amssymb% 


,russian%               remove in Eng. 


,izvru%                 Set izven% in Eng. 


% ,izvru-ks             for brief communications 


]{amsart} 


 


 


\newcommand{\tts}[1]{}


\newcommand{\rst}{\restriction}


\newcommand{\mer}{\leqslant}


\newcommand{\bor}{\geqslant}


\newcommand{\sk}[1]{{\bold #1^\prime}}


\newcommand{\dg}[1]{{\bold #1}}


\newcommand{\pp}{$d$-р.п.\ }


\newcommand{\op}{\oplus}


\newcommand{\dsum}{(D_1\op D_2)}


\newcommand{\bsum}{(B_1\op D_2)}


\newcommand{\ssum}{((D_1-D_2)\op A)}


\newcommand{\draz}{(D_1-D_2)}


\newcommand{\braz}{(B_1-B_2)}


\newcommand{\psum}{(\dsum\op A)}


\newcommand{\req}[2]{{\cal #1}_#2}


\newcommand{\con}[2]{#1\,\widehat{\ }\,#2}


\newcommand{\dom}{\operatorname{dom}}


\newcommand{\str}[1]{${\cal #1}$--стратеги}


\newcommand{\piv}[2]{$\langle\dg{#1},\dg{#2}\rangle$}


\newcommand{\limi}{\operatornamewithlimits{lim\,inf}}


 


\theoremstyle{definition} 


\theorembodyfont{\rm} 


\newtheorem{defnnonum}{\hskip\parindent Определение} 


\newtheorem{sled}{\hskip\parindent Следствие}


\renewcommand{\thedefnnonum}{} 





\theoremstyle{plain} 


\theorembodyfont{\it} 


\newtheorem{lemm}[thm]{\hskip\parindent Лемма}


\newtheorem{lth}{\hskip\parindent Лемма}


% was \newtheorem{lth}{\hskip\parindent Лемма}[thm]





\renewcommand{\theenumii}{\asbuk{enumii}}


%% use \alph in Eng.


\renewcommand{\labelenumi}{(\theenumi)}


\renewcommand{\labelenumii}{(\theenumii)}


\renewcommand{\thelth}{\thethm.\arabic{lth}}





\newcounter{n}


\setcounter{n}{1}





%\hoffset=-15mm%                  For Eng. 


 


\setcounter{page}{20} 


\god{1998} 


\nomer{2~(429)} %                Remove in Eng. 


%\nomer{Vol.42, No.2}%        Set in Eng. 


% \str{75--81}%                    Set in Eng. 


\udk{510.57} 


\author{\it А.А.\,Ефремов} 


% NB:   ^^ remove in Eng.


\title{Изолированные сверху $d$-р.п. степени, I} 


 


\thanks{Работа выполнена при частичной финансовой


поддержке Российского фонда фундаментальных исследований 


(номер проектов 93-011-16004, 96-01-00830), Международного научного


фонда Сороса (грант NNL000) и гранта Новосибирского университета.}





\date{} 


% \pagestyle{myheadings}%         Set in Eng. 


 


\pagestyle{plain} %              Remove in Eng. 


\begin{document} 


 


% \thispagestyle{plain}%          Set in Eng. 


 


\maketitle 


 





\section{Введение}


Множество  $A\subseteq\omega$ называется $d$--{\it рекурсивно перечислимым}


(\pp), если существуют рекурсивно перечислимые (р.п.) множества  $A_1$ и


$A_2$ такие, что  $A=A_1-A_2$. Тьюрингова степень называется \pp степенью,


если она содержит \pp множество; \pp степень называется {\it собственно} \pp, 


если она не является р.п. степенью (не содержит р.п. множеств).





В данной статье изучается взаимодействие р.п. и \pp степеней. С.Б.\,Купер и 


К.\,Йи \cite{CoYi} ввели определение {\it изолированой} \pp степени следующим


образом: \pp степень $\dg{d}$ называется изолированной, если существует р.п.


степень $\dg{a}<\dg{d}$ такая, что между ними нет р.п. степеней.


Естественным образом можно расширить определение изолированных \pp степеней,


например, изучать не нижний конус р.п. степеней под \pp степенью, а верхний


конус р.п. степеней над \pp степенью. Тогда изолированную степень в терминах 


предыдущего определения можно назвать {\it изолированной снизу} (ИСН) и по 


аналогии ввести новое





\begin{defnnonum}


\pp степень $\dg{d}$ называется {\em изолированной сверху}


({\em ИСВ\,}), если существует р.п. степень $\dg{b}>\dg{d}$ такая, что


между $\dg{d}$ и $\dg{b}$ нет р.п. степеней. В противном случае она называется


неизолированной сверху. Будем говорить, что степени $\dg{d}$ и $\dg{b}$


образуют {\em изолированную сверху пару}, и обозначать \piv{d}{b}.


\end{defnnonum}





В 1984 г. Л.\,Хэй и Р.\,Шор (неопубликовано) прямой конструкцией 


построили \pp степень 


$\dg{d}$ такую, что между $\dg{d}$


и $\sk{0}$ нет р.п. степеней, т.е. фактически


доказали существование изолированной сверху пары \piv{d}{0'}.


После того, как С.Б.\,Купер, С.\,Лемпп и П.\,Вотсон \cite{CLW}


доказали, что \pp степени 


плотны в структуре р.п. степеней, естественно ставить вопрос о плотности в р.п.


степенях любого класса \pp степеней с определенными свойствами. Отметим, что


вопросы о плотности ИСН \pp степеней и неизолированных снизу \pp степеней, 


поставленные в \cite{CoYi}, решены положительно в \cite{LF} и


\cite{ALS} соответственно.


Относительно ИСВ \pp степеней в этом направлении мы показываем, что их


``достаточно много''. А именно, в теореме~\ref{iso} утверждается,


что над любой 


неполной р.п. степенью в том же самом классе скачков существует ИСВ пара. 


Более того, оказалось, что свойство плотности для ИСВ \pp степеней в р.п. 


степенях не выполняется, т.к. существует нерекурсивная р.п. степень,


под которой


все \pp степени являются неизолированными сверху. Доказательству этого 


утверждения будет посвящена вторая часть статьи. Необходимо отметить, что


ИСВ пара в теореме~\ref{iso} строится в виде $\langle\text{dg\,}\draz,


\text{dg\,}\dsum\rangle$, и доказывается утверждение, с помощью 


которого {\it все} ИСВ пары можно представить в таком виде.





\section{Обозначения и терминология}





Обозначения стандартны и в основном мы следуем \cite{So}.


Рассматриваются множества и функции на натуральных числах $\omega=\{0,1,


2,3,\dotsc\}$. Обычно греческими буквами $\Phi$, $\Psi,\dotsc$, $\varphi$,


$\psi,\dotsc$ будут обозначаться частично рекурсивные (ч.р.) функции;


$f$, $g$, $h,\dotsc$ --- общерекурсивные (о.р.) функции; 


прописными латинскими буквами $A$, $B$, $C,\dotsc$ --- подмножества $\omega$.


Для ч.р. функции $\Phi$ символ $\Phi(x)\downarrow$ означает, что $x\in\dom


\Phi$; в противном случае пишется $\Phi(x)\uparrow$. Множество $A$


отождествляется с его характеристической функцией $\chi_A$. Запись $f\rst x$


означает ограничение $f$ на аргументы, меньшие чем $x$; аналогично для 


множеств; $A\subset B$ означает, что $A\subseteq B$, но $A\ne B$;


$A\op B=\{2x\mid x\in A\}\cup\{2x+1\mid x\in B\}$;


$\Phi_e$, $W_e$ ($\Phi_e^X$, $W_e^X$) означает $e$-ю


ч.р. функцию и ее область определения (с оракулом $X$) в некоторой 


фиксированной стандартной нумерации; наибольшее число, используемое в 


вычислении


функционалов $\Phi^X$, $\Psi^X,\dotsc$ в точке $x$, будет обозначаться 


соответствующей строчной греческой буквой $\varphi(x)$, $\psi(x)$.


Таким образом, изменение $X\rst(\gamma(x)+1)$ будет позволять изменять


значение $\Gamma^X(x)$. Запись $\mer_{\mathrm T}$ означает Тьюрингову


сводимость, а $\equiv_{\mathrm T}$ --- отношение эквивалентности, индуцируемое


ею. 





Когда речь идет о деревьях, $\alpha$, $\beta$, $\gamma,\dotsc$ будут 


означать вершины дерева, $|\alpha|$ --- длину~$\alpha$, $\con


{\alpha}{\beta}$ --- конкатенацию $\alpha$ и $\beta$; $\alpha


\subseteq\beta$ $(\alpha\subset\beta)$ означает, что $\alpha$ является


(собственным) начальным сегментом $\beta$; $\alpha<_{\mathrm L}\beta$ 


означает, что для некоторого $n\;$ ${\alpha\rst n}={\beta\rst n}$ и $\alpha(n)


<_{\Lambda}\beta(n)$ (где $<_{\Lambda}$ --- соответствующий порядок на дереве,


возможно различный на разных уровнях и $T\subseteq\Lambda^{<\omega}$);


$\alpha\mer\beta$ $(\alpha<\beta)$ означает, что $\alpha<_{\mathrm L}\beta$


или $\alpha\subseteq\beta$ $(\alpha\subset\beta)$. Множество $[T]$ 


бесконечных путей через дерево $T\subseteq\Lambda^{<\omega}$ определяется


как $\{h\in\Lambda^{\omega}\mid(\forall n)[{h\rst n}\in T]\}$.


 


Теорема будет доказываться методом приоритета с бесконечными нарушениями


с использованием дерева стратегий. 


План доказательства таков: 1)~выписывается бесконечный список


требований, которые нужно удовлетворить (выполнить); 2)~для каждого вида 


требований приводится стратегия, которую называем {\it основным модулем}


для удовлетворения требований данного вида; 3)~рассматриваются возможные


{\it выходы} (результаты работы) основных модулей для всех видов требований;


4)~определяется дерево стратегий и каждой вершине дерева назначается некоторое 


требование; 5)~описывается основная конструкция, которая определяет, как 


работают


и взаимодействуют стратегии на дереве; 6)~доказываются леммы,


показывающие, что построенные в процессе конструкции объекты удовлетворяют


условиям теоремы. Более подробно метод приоритета, использующий деревья,


изложен в (\cite{So}, гл.\,14; \cite{CLW} для \pp степеней).





При описании стратегий используем следующую терминологию: при выборе 


представителя


цикла слова ``достаточно большое число'' означают первое число, большее всех


упоминаемых в конструкции к данному моменту; {\it начинаем\/} цикл, 


позволяя ему


выполнить (1) и перейти в (2); {\it останавливаем} цикл, прекращая его работу,


определяя его запрет равным нулю и заставляя перейти в (1); {\it разрушаем}


цикл, останавливая его и считая, что его часть функционала становится 


неопределенной; {\it инициализируем} стратегию, разрушая все ее циклы и 


начиная


цикл 0; цикл {\it работает}, переходя из некоторого (кроме (1)) состояния в


другое; стратегия {\it работает}, позволяя циклу с наименьшим номером, 


который может это сделать, работать (в противном случае ничего не делает).





В процессе наших конструкций будем строить р.п., \pp множества и ч.р.


функционалы по шагам. Для обозначения полученного


результата к концу шага $s$ для строящихся и данных объектов будем 


использовать запись типа $A[s]$, $\Phi_e[s]$, $A=\Phi_e [s]$, но, когда


это не ведет к двусмысленности, параметр $[s]$ будем опускать.


При необходимости


подчеркнуть, что объекты рассматриваются как результат различных шагов,


будем использовать запись типа $A_s$, $\varphi_{e,t}$.


Параметры, получившие некоторое значение в ходе конструкции, остаются 


с этим значением, пока не будут переопределены. 





\section{Основной результат}


 


\begin{thm}\label{iso}


Над любой р.п. степенью $\dg{a}<\sk{0}$ существует изолированная


сверху пара $\langle\dg{d},\dg{b}\rangle$ {\em (}т.е. $\dg{a}<


\dg{d}<\dg{b}$ {\em )} такая, что $\sk{a}=\sk{b}$.


\end{thm}





{\bf Доказательство.} Пусть $A\in\dg{a}$ --- данное р.п. множество.


Будем строить р.п. множества $D_1,D_2\subseteq D_1$ такие, что


$\deg((D_1-D_2)\op A)$ и $\deg\psum$ будут


образовывать искомую пару. Для этого достаточно удовлетворить


для всех $e$ следующим требованиям:


\begin{itemize}


\item[$\req{P}{e}:$] $\dsum\ne\Phi_e^{\ssum}$,\\


где $\{\Phi_e\}_{e\in\omega}$ --- перечисление ч.р. функционалов;


\item[$\req{R}{e}:$] $\draz=\Phi_e^{W_e}\Longrightarrow


(\exists\Delta_e)(\dsum=\Delta_e^{W_e})$,\\


где $\{\Phi_e,W_e\}_{e\in\omega}$ --- перечисление пар, состоящих из 


ч.р. функционала


$\Phi_e$ и р.п. множества~$W_e$;


\item[${\cal N}:$] $\psum^\prime\mer_{\mathrm T}A^\prime$.


\end{itemize}


Необходимо отметить, что $\ssum\rst x$ и $\psum\rst x$ в оракуле


рассматривается как $(\draz\rst x)\op(A\rst x)$ и 


$(\dsum\rst x)\op(A\rst x)$ cоответственно.


\medskip





Рассмотрим основные модули для удовлетворения этих требований.


\medskip





\underline{\bf Основной модуль для $\req{P}{e}:$}


$\dsum\ne\Phi_e^{\ssum}$


\medskip





Стратегия работает по циклам, каждый из которых действует следующим


образом.


\begin{enumerate}


\item[] \underline{\bf Цикл $k$}


\item Выбираем представителя цикла $x_k$ достаточно большим.


\item Ждем шага $s_1: \Phi_e^{\ssum}(2x_k)\downarrow=\Phi_e^{\ssum}


(2x_k+1)\downarrow=0 [s_1]$.


\item Определяем функционал $\Gamma_e^A(k)[s_1+1]=K(k)[s_1]$ с


$\gamma_e(k)[s_1+1]=\varphi_e(2x_k+1)[s_1]$, где $K$ --- креативное


множество. Определяем запрет $r(k)[s_1+1]=\varphi_e(2x_k+1)


[s_1]$. Начинаем $(k+1)$-й цикл. Одновременно продолжаем работу


нашего цикла, переходя в~(4).


\item Ждем изменения $A\rst\gamma_e(k)$ или $K(k)$.


\begin{enumerate}


\item Изменение $A\rst\gamma_e(k)$ произошло раньше. Тогда


разрушаем циклы с номерами $k'>k$, полагаем запрет $r(k)=0$,


$\Gamma_e^A(k)$ считаем неопределенным (из-за изменения $A\rst


\gamma_e(k)$) и идем в~(2).


\item Изменение $K(k)$ произошло раньше. Тогда останавливаем


циклы с номерами $k'>k$, запрет $r(k)=\varphi_e(2x_k+1)$


сохраняется, идем в~(5).


\end{enumerate}


\item Перечисляем $x_k$ в $D_1$ и в $D_2$ (тем самым


$\draz(x_k)$ не изменилось, а $\dsum$ в точках $2x_k$ и


$2x_k+1$ стало равным~1).


\item Ждем шага $t$, на котором произойдет изменение $A\rst\gamma_e(k)$.


\item Разрушаем циклы с номерами $k'>k$, переопределяем функционал


$\Gamma_e^A(k)[t+1]=K(k)[t]$ с $\gamma_e(k)[t+1]=\varphi_e(2x_k+1)+1$,


полагаем $r(k)[t+1]=0$, начинаем $(k+1)$-й цикл и идем в~(8).


\item Конец работы.


\end{enumerate}





\underline{\bf Выходы $\req{P}{e}$-стратегии}


\begin{itemize}


\item[1)] Существует шаг $s_0$, после которого ни один цикл не работает.


Тогда некоторый цикл $k_0$ остается навсегда в~(2) или в~(6).


Требование удовлетворено. Общий запрет всех циклов имеет конечный предел.


Обозначим этот выход через $f$ (конечный).


\item[2)] Работает бесконечное число циклов, но существует цикл с наименьшим


номером $k_0$, который бесконечно много раз переходит из~(4) (через


(a)) в~(2). Требование удовлетворено, т.к. функционал 


$\Phi_e^{\ssum}$ в точке $2x_{k_0}+1$ расходится. Общий запрет всех


циклов $R(e)$ может быть бесконечным в пределе, но существует


конечный $\limi\limits_s R(e)$. Номер цикла $k_0$ и будет обозначать


этот выход.


\item[3)] Работает бесконечное число циклов, но каждый работает конечное число


раз, останавливаясь в~(4) или в~(8). Требование не выполнено и построен


всюду определенный функционал $\Gamma_e^A$, который правильно


вычисляет креативное множество $K$ в каждой точке. Так как


$A<_{\mathrm T}K$ по условию теоремы, то этот выход невозможен в нашей


конструкции, и его учитывать не будем.


\end{itemize}





\underline{\bf Основной модуль для $\req{R}{e}:$}


$\draz=\Phi_e^{W_e}\Longrightarrow(\exists\Delta_e)


(\dsum=\Delta_e^{W_e})$


\medskip





Стратегия работает по циклам, каждый из которых действует следующим образом.


\pagebreak


\begin{enumerate}


\item[] \underline{\bf Цикл $k$}


\item Выбираем представителя цикла $x_k$ достаточно большим.


\item Ждем шага $s: \draz\rst(x_k+1)=\Phi_e^{W_e}\rst(x_k+1) [s]$.


\item (Пере)определяем функционал $\Delta_e^{W_e}(k)[s+1]=\dsum(c) [s]$


с $\delta_e(k)[s+1]=\varphi_e(x_k)[s]$. Начинаем $(k+1)$-й


цикл. Одновременно продолжаем работу нашего цикла, переходя в~(4).


\item Ждем изменения $W_e\rst\delta_e(k)$ или $\dsum(k)$.


\begin{enumerate}


\item Изменение $W_e\rst\delta_e(k)$ произошло раньше. Тогда


разрушаем циклы с номерами $k'>k$, перечисляем $x_{k'}$ в


$D_2$, если $x_{k'}\in D_1$, $\Delta_e^{W_e}(k)$ считаем


неопределенным и идем в~(2).


\item Изменение $\dsum(k)$ произошло раньше. Тогда останавливаем


циклы с номерами $k'>k$, идем в~(5).


\end{enumerate}


\item Перечисляем $x_k$ в $D_1$.


\item Ждем шага $t: \draz\rst(x_k+1)= \Phi_e^{W_e}\rst(x_k+1) [t]$.


\item Переопределяем функционал $\Delta_e^{W_e}(k)[t+1]=\dsum(k) [t]$.


Перечисляем $x_k$ в $D_2$, разрушаем циклы с номерами


$k'>k$, перечисляем $x_{k'}$ в $D_2$, если $x_{k'}\in D_1$,


и идем в~(2).


\end{enumerate}





\underline{\bf Выходы $\req{R}{e}$--стратегии}


\begin{itemize}


\item[1)] Работает конечное число циклов. В этом случае некоторый цикл $k_0$


либо ждет в~(2) или в~(6), либо бесконечное число раз переходит из~(4) в~(2).


Требование $\req{R}{e}$ удовлетворено, т.к. $\draz(x_{k_0})\neq\Phi_e^


{W_e}(x_{k_0})$. Функционал $\Delta_e^{W_e}$ не является всюду 


определенным в этом случае. Отметим, что перечисляется только конечное число


элементов в~$D_1$ и~$D_2$. Обозначим этот выход через~1.


\item[2)] Работает бесконечное число циклов. Возможны два случая.


 


(а) (настоящий бесконечный выход) Для каждого $k$ существует шаг $s_k$,


после которого цикл с номером $k$ не работает. Тогда каждый из


циклов окончательно ждет в (4). Требование $\req{R}{e}$ удовлетворено,


т.к. в этом случае функционал $\Delta_e^{W_e}$ всюду определенный и 


правильно вычисляет $\dsum$.





(б) (мнимый бесконечный выход) Существует цикл $k_0$, который бесконечное


число раз переходит из~(4) в~(2). Требование $\req{R}{e}$ удовлетворено, 


т.к. $\Phi_e^{W_e}(x_{k_0})\uparrow$. Функционал $\Delta^{W_e}$ в этом 


случае не является всюду определенным, т.к. мы его постоянно в точках 


$k\bor k_0$ разрушаем и $\delta_e(k)\to\infty$.





$\qquad$Заметим, что в обоих случаях все элементы, положенные


в~$D_1$, затем


перечисляются в~$D_2$ и на $\draz$ не оказывается никакого влияния, но


в $\dsum$ кладется бесконечно много чисел. Возникает проблема между 


требованием $\req{R}{e}$ и требованием $\req{N}{j}$, которое 


расположено в дереве стратегий ниже бесконечного выхода $\req{R}{e}$. 


Эта проблема решается 


в процессе конструкции и ей посвящено доказательство леммы~\ref{R-N}.


Обозначим этот выход через~0.


\end{itemize}





\underline{\bf Требование ${\cal N}:$}


$\psum'\mer_{\mathrm T}A'$


\medskip





Наша стратегия будет естественной релятивизацией известной стратегии


построения низкого множества (\cite{So}, с\,.111). Таким образом, будем


строить множество $\psum$ низким относительно $A$ (аналогичная


идея использовалась в~\cite{LSl}).Для этого требование ${\cal N}$ 


разобьем на бесконечный список требований


$$


\req{N}{e}: (\exists^\infty s)[(\Phi_{e}^{\psum}(e)


\downarrow[s])\, \& \,


(A_s\rst\varphi_{e,s}(e)=A\rst\varphi_{e,s}(e))]


\Longrightarrow\Phi_e^{\psum}(e)\downarrow


$$


для всех $e\in\omega$, где $\{\Phi_e\}_{e\in\omega}$ ---


перечисление ч.р. функционалов.





Покажем, как при условии удовлетворения требований $\req{N}{e}$,


$e\in \omega$, выполняется требование ${\cal N}$.


Определим $A$-рекурсивную функцию


$$


g^A(e,s)=


\begin{cases}


1, & \text{если }\Phi_{e}^{\psum}(e)\downarrow[s]\,\&\,


A_s\rst\varphi_{e,s}(e)=A\rst\varphi_{e,s}(e);\\


0 & \text{в противном случае}.


\end{cases}


$$


Если требования $\req{N}{e}$ для всех $e\in\omega$ удовлетворены,


то $\Hat g(e)=\lim\limits_s g^A(e,s)$ существует для всех $e$.


По лемме о пределе (\cite{So}, с.\,57) $\Hat g\mer_{\mathrm T}A'$, и, т.к.


$\Hat g$ является характеристической функцией $\psum'$, имеем


$\psum'\mer_{\mathrm T}A'$.


\medskip





\underline{\bf Основной модуль для $\req{N}{e}$}


\medskip





Стратегия для удовлетворения требования $\req{N}{e}$ работает


следующим образом.


\begin{enumerate}


\item Ждем шага $s_1: \Phi_e^{\psum}(e)\downarrow [s_1]$.


\item Устанавливаем запрет $r(e)[s_1+1]=\varphi_{e}(e)[s_1]$ на $\dsum$.


\item Ждем шага $s_2$, на котором произошло изменение $A\rst\varphi_{e}(e)$.


\item Устанавливаем запрет $r(e)[s_2+1]=0$, и идем в~(1).


\end{enumerate}





\underline{\bf Выходы $\req{N}{e}$-стратегии}


\begin{itemize}


\item[1)] Существует шаг $s$, после которого стратегия не работает. Тогда


она ждет либо в~(1), либо в~(3). В этом случае запрет $r(e)$ имеет


конечный предел. Обозначим этот выход через~1.


\item[2)] Стратегия работает бесконечно много раз, переходя из~(4) в~(1).


В этом случае существует $\limi\limits_s r(e)=0$. Обозначим этот выход


через~0.


\end{itemize}





\underline{\bf Дерево стратегий}





Обозначим множества выходов стратегий ${\cal P}$, ${\cal R}$,


${\cal N}$ через $\Lambda_P$, $\Lambda_R$ и $\Lambda_N$ соответственно,


т.е.


$\Lambda_P=


\{{\it 0}<_{\Lambda_P}{\it 1}<_{\Lambda_P}{\it 2}<_{\Lambda_P}\dotsb


<_{\Lambda_P}f\,\}$, $\Lambda_R=\{{\it 0}<_{\Lambda_R}{\it 1}\,\}$,


$\Lambda_N=\{{\it 0}<_{\Lambda_N}{\it 1}\,\}$. Пусть $\Lambda=\Lambda_P


\cup\Lambda_R\cup\Lambda_N$. Тогда деревом стратегий назовем множество


$$


T=\{\alpha\in\Lambda^{<\omega}\bigl.\bigr|


(\forall k<|\alpha|)[\alpha(k)\in\Lambda_P,\;


\Lambda_R,\; \Lambda_N \text{ для } k\equiv 0,1,2\pmod 3]\}.


$$





Каждой вершине $\alpha\in T$ назначаем некоторое требование, над которым


она будет работать. Пусть $|\alpha|=3e+k$,


где $e,k\in\omega$ и $k\mer 2$. Тогда


$$


\text{назначаем вершине }\alpha \text{ требование }


\begin{cases}


{\req{P}{e}}, &\text{если $k=0$};\\


{\req{R}{e}}, &\text{если $k=1$};\\


{\req{N}{e}}, &\text{если $k=2$},


\end{cases}


$$


т.е. каждый уровень дерева работает над одним требованием.


 


Теперь $\alpha$-стратегией будем называть вариант основного модуля для 


требования, которое назначено вершине $\alpha$, с некоторыми изменениями,


описанными в конструкции; $S$-стратегией, где


$S\in\{{\cal P},{\cal R},{\cal N}\}$, будем называть стратегию 


$\alpha$, если она работает над требованием вида $\req{P}{e}$, 


$\req{R}{e}$, $\req{N}{e}$ соответственно. Шаг $s$ назовем $\alpha$-шагом, 


если стратегия $\alpha$ имела возможность работать на шаге $s$.





\bigskip


\underline{\bf Полная конструкция}


\medskip





\noindent\underline{\bf Шаг 0.} $D_{1,0}= D_{2,0}=\emptyset$, все $\Delta_e$,


$\Gamma_e$ неопределены, все стратегии $\alpha\in T$


инициализированы.





\noindent\underline{\bf Шаг $s+1$} состоит из двух фаз.





1-я фаза. Ищем самую левую \str{P}ю $\alpha\in T$ такую, что


некоторый цикл с наименьшим номером $k_0$ ждет в~(4) и 


$K(k_0)[s]\ne K(k_0)[s_1]$, где $s_1$ --- шаг, на котором 


цикл $k_0$ стратегии $\alpha$ прошел через~(2) в~(3). Если такая 


найдется, то позволяем этому циклу перейти через~(4б) в~(5) и выполнить


его. Затем инициализируем все стратегии $\beta\bor\con{\alpha}{k_0}$.


Если такой стратегии нет, то ничего не делаем.





2-я фаза. Работаем по подшагам $t\mer 3s+2$,


причем на каждом подшаге $t$ одна из стратегий $\alpha$ длины


$t$ будет иметь возможность работать. Возможны три случая.


 


$|\alpha|=t=3e$. Стратегия $\alpha$ работает над требованием


${\cal P}_e$, как описано в основном модуле для ${\cal P}_e$.


Определяем выход $o$ стратегии $\alpha$ следующим образом:


если некоторый цикл работал, то выходом будет номер этого цикла, в


противном случае --- $f$. Далее, инициализируем стратегии


$\beta>_{\mathrm L}\con{\alpha}{o}$ и определяем следующую


стратегию $\beta$, которая имеет возможность работать на подшаге 


$t+1$: если работающий цикл $k_0$ перешел из~(4) в~(2), то 


$\beta=\con{\alpha}{k_0}$, в противном случае 


$\beta=\con{\alpha}{f}$. Переходим к подшагу $t+1$.





$|\alpha|=t=3e+1$. Стратегия $\alpha$ работает над требованием


${\cal R}_e$, как описано в основном модуле для ${\cal R}_e$ с


небольшими изменениями, а именно: если между двумя последовательными


$\alpha$-шагами $t_1<t_2$ некоторый цикл $k_0$ ждет в~(4), и


изменилось сначала значение $\dsum(k_0)$, а затем $W_e\rst


\delta_e(k_0)$, то на шаге $t_2$ цикл $k_0$ (если ему дадут


работать) переходит из~(4) в~(2) через~(4а) (т.е. не смотрит на то, что


изменение $\dsum(k_0)$ произошло раньше).


Определяем выход $o$ стратегии $\alpha$ следующим образом:


если некоторый цикл работал впервые, то выходом считаем 0, в


противном случае --- 1. Далее, инициализируем стратегии


$\beta>_{\mathrm L}\con{\alpha}{o}$. 


Если работающий цикл прошел через~(5) или~(7), то


определяем $\delta_{s+1}=\alpha$, инициализируем стратегии


$\beta\supset\con{\alpha}{{\it 1}}$ и переходим к следующему шагу


$s+2$. В противном случае определяем следующую стратегию $\beta$,


которая имеет возможность работать на подшаге $t+1:\beta=\con{\alpha}{o}$.


Переходим к подшагу $t+1$.


 


$|\alpha|=t=3e+2$. Стратегия $\alpha$ работает над требованием


${\cal N}_e$, как описано в основном модуле для ${\cal N}_e$.


Определяем выход $o$ стратегии $\alpha$ следующим образом:


если произошел переход из~(3) в~(1), то выход --- 0, в


противном случае --- 1. Далее, инициализируем стратегии


$\beta>_{\mathrm L}\con{\alpha}{o}$. Если был выполнен~(2), то


инициализируем еще стратегии $\beta\supset\con{\alpha}{o}$.


Если $t=3s+2$, то определяем $\delta_{s+1}=\alpha$ и переходим


к следующему шагу $s+2$.


В противном случае определяем стратегию $\beta=\con{\alpha}{o}$,


которая имеет возможность работать на подшаге $t+1$.


Переходим к подшагу $t+1$.


Конструкция закончена.





\bigskip


\underline{\bf Проверка конструкции}


\medskip


 


Для завершения доказательства теоремы докажем ряд лемм.


 


Определим истинный путь $f\in [T]$ как самую левую бесконечную ветвь


дерева $T$, посещаемую $\delta_s$ в течение конструкции бесконечное 


число раз, т.е. для всех $n$, если $\alpha=f\rst n$, то $f(n)=a$


означает окончательный выход стратегии $\alpha$. Покажем, что истинный 


путь в нашей конструкции определяется корректно. 





\begin{lth}


Истинный путь $f$ существует.


\end{lth}





\begin{pf}


Доказательство ведем индукцией по глубине дерева $n$. Полагаем 


$f\rst 0=\lambda$. Пусть $\alpha=f\rst n$ уже определено. Покажем,


что $f(n)$ определяется единственным образом. Возможны три случая.





1) $|\alpha|=n=3e$, т.е. $\alpha$ является \str{P}ей. В этом случае 


$f(n)=\limi\limits_s \delta_s(n)$ в том смысле,


что $(\exists^{<\infty}s)[\delta_s(n)<_{\mathrm L}f(n)]$ и


$(\exists^{\infty}s)[f(n)=\delta_s(n)]$.


То есть либо $f(n)=k_0$, где $k_0$ --- номер


цикла, бесконечно много раз переходящего из~(4) в~(2), либо $f(n)=f$. 


Заметим, что из предположения $f=\alpha$ получаем $K\mer_{\mathrm T}A$,


что противоречит условию теоремы. 


 


2) $|\alpha|=n=3e+1$, т.е. $\alpha$ является \str{R}ей.


Определяем $f(n)=\limi\limits_s\delta_s(n)$, т.к. в этом


случае предел существует всегда. Легко показывается, что в случае, когда


$f(n)=0$, $(\exists^{\infty}s)[\con{\alpha}{\it 0}\subseteq \delta_s]$,


т.е. конструкция ``не застревает'' на \str{R}и.





3) $|\alpha|=n=3e+2$, т.е. $\alpha$ является \str{N}ей. Определяем 


$f(n)=\limi\limits_s\delta_s(n)$, т.к. в этом случае предел существует


всегда.


\end{pf}


 


\begin{lth}\label{INI}


Все вершины $\alpha\subset f$ инициализируются конечное число раз.


\end{lth}





\begin{pf}


Вначале покажем, что из-за 1-й фазы конструкции $\alpha$ инициализируется


конечное число раз. Действительно, т.к. $\alpha\subset f$, то стратегии


$\gamma<_{\mathrm L}\alpha$ в первой фазе могут работать конечное число раз.


Что касается \str{P}й $\con{\gamma}{k}\subseteq\alpha$, то любой 


цикл $k\;$ \str{P}и может перечислить своих представителей лишь один раз.


Так как


таких $\gamma$ лишь конечное число, то $\alpha$ в этом случае


инициализируется конечное число раз.





Из-за 2-й фазы конструкции $\alpha$ может инициализироваться


стратегиями $\beta$, которые либо $\beta<_{\mathrm L}\alpha$, либо


$\beta\subset\alpha$. В первом случае, т.к. $\alpha\subset f$, 


$\beta$ может это делать только конечное число раз. Во втором случае,


если $\alpha$ находится под конечным выходом $\beta$, то $\alpha$


инициализируется конечное число раз; в противном случае $\beta$ вообще


не инициализирует $\alpha$. 


\end{pf}





\begin{lth}\label{R-N}


Запрет \str{N}и $\alpha\subset f$ нарушается конечное число раз. 


\end{lth}





\begin{pf}


Единственным заслуживающим внимания является положение, когда \str{N}я


$\alpha\subset f$ расположена под бесконечными выходами \str{R}й. 


Достаточно рассмотреть два случая.


 


1) Пусть \str{R}я одна: $\con{\beta}{{\it 0}}\subseteq\alpha


\subset f$. Первоначально \str{R}я $\beta$ может перечислить некоторый 


элемент $x$ в $\dsum$ только из-за перечисления \str{P}ей $\gamma$


в первой фазе конструкции некоторого элемента $y<x$. Самым интересным 


является случай, когда $\con{\beta}{{\it 0}}\subseteq\gamma\subset\alpha$, 


т.к. в этом случае $\alpha$ инициализируется, а $\beta$ нет. Однако


$\alpha$ будет иметь возможность работать (устанавливать свой


новый запрет) только после того, как $\beta$ пройдет через~(5), (6), (7) 


и~(2), т.е. $\beta$ ничего не нарушает.





2) Пусть имеются две \str{R}и: $\con{\beta_1}{{\it 0}}\subset\con


{\beta_2}{{\it 0}}\subseteq\alpha$, и некоторая \str{P}я $\gamma$ в


первой фазе конструкции перечислила некоторый элемент~$x$ в~$\dsum$. Опять


самый интересный случай, когда $\con{\beta_1}{{\it 0}}\subset\con


{\beta_2}{{\it 0}}\subseteq\gamma\subset\alpha$. Обозначим через $y_1(k)$ и


$y_2(k)$ представителей $k$-го цикла стратегий $\beta_1$ и $\beta_2$


соответственно. Рассмотрим ситуацию, когда $\beta_2$ создает работу для


$\beta_1$. Это может происходить следующим образом: $\beta_1$ выполняет


свой цикл $x$ и, корректируя свой функционал в точке $x$, перечисляет~$y_1(x)$


вначале в~$D_1$, а затем в~$D_2$; до того, как $\beta_2$ получит


возможность работать, $\beta_1$ успевает запустить свой $y_2(x)$-й


цикл; теперь $\beta_2$ корректирует свой функционал в точке $x$,


перечисляя~$y_2(x)$ в~$D_1$; после этого $\beta_1$ вынужден


корректировать функционал в точке~$y_2(x)$, перечисляя~$y_1(y_2(x))$.


На этом процесс взаимной работы $\beta_1$ и $\beta_2$ заканчивается.


Не вдаваясь в подробности, отметим, что $\beta_1$ для $\beta_2$ 


работы не создает. Но заметим, что по конструкции до $\alpha$ мы доберемся 


только после того, как вся эта работа будет завершена.


 


Если рассмотреть любую конечную последовательность $\con{\beta_1}


{{\it 0}}\subset\con{\beta_2}{{\it 0}}\subset\dotsc\subset\con{\beta_n}


{{\it 0}}\subseteq\alpha$, то совершенно аналогично можно показать, что


процесс взаимной работы $\beta_1,\dotsc,\beta_n$ конечен и никак не


затрагивает запрет стратегии $\alpha$.





Делаем вывод, что нарушения запрета $\alpha$ происходят только из-за


перечисления элементов в~$\dsum$ в первой фазе конструкции \str{P}ями


$\gamma$ такими, что $\gamma<\alpha$. Это всегда сопровождается


инициализацией, а по лемме~\ref{INI} это происходит конечное число раз.


\end{pf}





Обозначим через $R(\alpha,s)$ запрет стратегии $\alpha$ после шага


$s$ (в случае \str{P}и это будет общий запрет всех циклов).





\begin{lth}\label{RES}


Любая стратегия $\alpha\subset f$ для $\gamma\supset\alpha$ на $f$


создает конечный запрет $R(\alpha)$.


\end{lth}





\begin{pf}


В нашей конструкции запреты устанавливают только ${\cal P}$- и \str{N}и.


Пусть $\alpha\subset f$ --- \str{P}я. Как отмечалось в разделе о выходах


\str{P}и, если $\con{\alpha}{{\it f\/}}\subset f$, то $R(\alpha)=\lim\limits_s


R(\alpha,s)$ существует и конечен. Если $\con{\alpha}{{\it k_0\/}}\subset f$,


то существует конечный $\limi\limits_s R(\alpha,s)=R(\alpha)$.





Пусть $\alpha\subset f$ --- \str{N}я. Как описано в разделе о выходах \str{N}и


в случае $\con{\alpha}{{\it 1}}\subset f$, $R(\alpha,s)$ имеет конечный


предел, а в случае $\con{\alpha}{{\it 0\/}}\subset f\;$


$\limi\limits_s R(\alpha,


s)=0$. Обозначим эти пределы через $R(\alpha)$.


\end{pf}





\begin{lth}\label{REQ}


Все требования вдоль истинного пути удовлетворяются.


\end{lth}





\begin{pf}


Отметим, что по построению дерева стратегий вдоль любой бесконечной ветви


расположены все требования.





Пусть $\alpha\subset f$ --- любая стратегия. По лемме~\ref{INI} существует шаг


$s$, после которого она уже не инициализируется. Используя лемму~\ref{RES}


и тот факт, что $\alpha\subset f$, легко показать, что $R=\max\{R


(\beta)\mid\beta<\alpha\}$ конечен. Так как установка нового запрета


сопровождается инициализацией право- и нижележащих стратегий, то все


вновь выбранные представители стратегии $\alpha$ после шага $s$ уже будут


больше $R$. На всех последующих $\alpha$-шагах наша стратегия


будет работать, как описано в ее основном модуле, и требование, над которым


она работает, будет выполнено.


\end{pf}





Доказательство теоремы закончено.


\setcounter{sled}{0}


\begin{sled}


Каждый класс скачков содержит бесконечную возрастающую последовательность


ИСВ \pp степеней.


\end{sled}





\begin{pf}


Пусть $\Bbb C$ --- любой класс скачков, кроме $\text{High}_1$,


и пусть р.п. степень


$\dg{b_0}\in\Bbb C$. Применяя теорему~1 с $\dg{a}=\dg{b_0}$, получаем ИСВ


\pp степень $\dg{d_0}$ и р.п. степень $\dg{b_1}$ такие, что $\dg{b_0}<\dg{d_0}<


\dg{b_1}$ и $\dg{b_1}\in\Bbb C$. Следовательно, и $\dg{d_0}\in\Bbb C$.


Применяя теорему~1 с $\dg{a}=\dg{b_1}$, получаем ИСВ \pp степень $\dg{d_1}$ и


р.п. степень $\dg{b_2}$ такие, что $\dg{d_0}<\dg{b_1}<\dg{d_1}<\dg{b_2}$ и


$\dg{b_2}\in\Bbb C$, и т.д.





Пусть теперь ${\Bbb C}=\text{High}_1$.


Добавим в формулировку теоремы~1 условие


$\dg{b}<\dg{0'}$. Для того чтобы доказать теорему в этом случае, 


заменим в списке


требований теоремы~1 требования вида ${\cal N}_e$ на требования


${\cal N\cal P}_e:


B\ne\Phi_e^{\psum}$, где $B$ --- вспомогательное р.п. множество, которое мы


строим. Эти требования гарантируют, что $\deg(\psum)\ne\dg{0'}$. Эти изменения


требований не добавляют в конструкцию принципиальных трудностей, и


доказательство остается почти в прежнем виде (с необходимыми изменениями в


дереве стратегий и полной конструкции). После этого, применяя теорему~1


с новым условием, получаем


$\dg{b_0}<\dg{d_0}<\dg{b_1}<\dg{0'}$ и $\dg{b_1}\in


\text{High}_1$, и т.д.


\end{pf}





Построенная в теореме 1 ИСВ пара имеет вид $\langle \deg\draz,


\deg\dsum\rangle$. Оказывается, что любую ИСВ пару можно представить


в таком виде. Вначале докажем следующее


\medskip





%% \begin{lemm}\label{tex}


{\bf Утверждение. }{\it Пусть даны произвольные \pp степень $\dg{d}$ и р.п.


степень $\dg{a}\bor\dg{d}$.


Тогда существует \pp множество $\draz\in\dg{d}$ такое, что $\deg{\dsum}\mer


\dg{a}$.}


%%\end{lemm}





\begin{pf} Пусть даны произвольное \pp множество $\braz\in\dg{d}$


и р.п. множество $A\in\dg{a}$. Построим \pp множество


$\draz=\braz$, удовлетворяющее условию $\dsum\mer_{\mathrm T}A$. Идея


состоит в том, чтобы при построении $\draz$ ``пропускать'' элементы, которые


входят в~$B_1$, а затем в~$B_2$ ``без разрешения $A$''. Опишем


нашу стратегию. 





Пусть $\braz=\Phi_e^A$. Строим $\draz=\braz$ и ч.р. функционал


$\Delta:\Delta^A=\dsum$. Наша стратегия работает по следующим циклам.


\begin{enumerate}


\item[] \underline{\bf Цикл $k$}


\item Ждем шага $s: \braz(k)=\Phi_e^A(k)\downarrow [s]$.


\item (Пере)определяем функционал $\Delta^A(2k)=\dsum(2k)=


(B_1\oplus B_2)(2k)$, $\Delta^A(2k+1)=\dsum(2k+1)=


(B_1\oplus B_2)(2k+1)$~$[s+1]$ c


$\delta(2k)=\delta(2k+1)=\varphi_{e,s}(k)$.


(Тем самым автоматически определяем $\draz(k)=\braz(k)$.)


Начинаем $(k+1)$-й цикл. Одновременно продолжаем работу нашего цикла,


переходя в~(3).


\item Ждем изменения $A\rst\delta(2k+1)$ или $\braz(k)$.


\begin{enumerate}


\item Изменение $A\rst\delta(2k+1)$ произошло раньше. Тогда


разрушаем циклы с номерами $k'>k$; $\Delta^A(2k),(2k+1)$ считаем


неопределенным (из-за изменения $A\rst\delta(k)$) и идем в~(1).


\item Изменение $\braz(k)$ произошло раньше. Тогда останавливаем


циклы с номерами $k'>k$, идем в~(4).


\end{enumerate}


\item Изменение могло произойти двумя путями:


\begin{enumerate}


\item $\braz(k)$ изменилось с~0 на~1. Тогда ждем восстановления


равенства на некотором шаге $t:\braz(k)=\Phi_e^A(k)\downarrow~[t]$.


Оно произойдет либо из-за изменения $\braz(k)$ на~0, либо


$A\rst\varphi_e(k)$. В первом случае продолжаем работу $(k+1)$-го


цикла и идем в~(5). Во втором случае перечисляем $k \mbox{ в } D_1 $,


разрушаем циклы с номерами $k'>k$ и идем в~(2) переопределять


$\Delta^A$ в точках $2k$ и $2k+1$.


\item $\braz(k)$ изменилось с~1 на~0. Тогда ждем восстановления


равенства на некотором шаге $t:\braz(k)=\Phi_e^A(k)\downarrow=0~[t]$.


Оно может произойти только из-за изменения


$A\rst\varphi_e(k)$. Перечисляем $k$ в $D_2$,


разрушаем циклы с номерами $k'>k$, переопределяем $\Delta^A(2k+1)=


\dsum(2k+1)$ с окончательным произвольным $\delta(2k+1)$,


начинаем $(k+1)$--й цикл и идем в~(5).


\end{enumerate}


\item Конец работы.


\end{enumerate}





Конструкция заключается в том, чтобы на каждом шаге $s>0$ стратегия


давала работу циклу с наименьшим номером, который имеет работу.





Для доказательства леммы отметим, что (т.к. $\braz=\Phi_e^A$)


окончательно все циклы будут находиться в~(3) или в~(5). Понятно, что


$\draz=\braz$и $\dsum\mer_{\mathrm T}A$.


\end{pf}





Из утверждения получаем





\begin{sled}\label{krit}


Собственно \pp степень $\dg{d}$ и р.п. степень $\dg{b}>\dg{d}$ образуют


изолированную сверху пару тогда и только тогда, когда существует \pp 


множество $\draz\in\dg{d}$ такое, что $\dsum\in\dg{b}$ и между


степенями


$\deg\draz$ и $\deg\dsum$ нет р.п. степеней.


\end{sled}





\begin{thebibliography}{10}





\bibitem{ALS} Arslanov M.M., Lempp S., Shore R.A.


{\em On isolating r.e. and 


isolated d-r.e. degrees} // London Math. Soc. -- Lecture Note Series.


-- 1996. -- V.\,224. -- P.\,61--80.





\bibitem{CLW} Cooper S.B., Lempp S., Watson P.


{\em Weak density and cupping


in the d-r.e. degrees} // Israel J. Math. -- 1989. -- V.67. -- \No\,2.





\bibitem{CoYi} Cooper S.B., Yi X., {\em Isolated d-r.e. degrees} //


Preprint. --  University of Leeds. -- 1995.





\bibitem{LF} LaForte G. {\em The isolated d.r.e. degrees are dense in the r.e.


degrees} // Math. Logic Quarterly. -- 1996. -- V.\,42. -- P.\,83--103.





\bibitem{LSl} Lempp S., Slaman T.A. {\em A limit on relative genericity


in the recursively enumerable sets} // J. Symb. Logic. -- 1989. -- V.\,54.


-- \No\,2.





\bibitem{So} Soare R.I. {\it Recursively enumerable sets and degrees}. --


Berlin: Springer-Verlag, 1987.





\end{thebibliography}





\vskip 2cm 


 


\post{\it Казанский государственный университет}{\it Поступила}


\post{}{07.06.1995} 


 


\end{document} 





